% Unit 098 terjemahan Bahasa Indonesia dari chapter7.tex baris 1968--2035.
% Sumber: Methods of Algebra, Volume 2, commit
% 9a5803ff77dd3257484cb177f851a73770a59dd3 (CC BY 4.0).
% stable-unit-id: o014.aljabr2.chapter7.morita-theory-c
% Cakupan: akhir Bagian 7.7, dari karakterisasi kategoris modul yang
% dibangkitkan secara hingga hingga bentuk simetris ekuivalensi Morita.

% segment-id: o014.aljabr2.chapter7.morita-theory-c.l001
\begin{lemma}\label{prop:fg-module-cat}
	Misalkan $R$ gelanggang dan $N$ modul-$R$ kanan. Maka $N$ dibangkitkan
	secara hingga jika dan hanya jika sifat berikut berlaku: untuk setiap
	keluarga subobjek $(N_i)_{i \in I}$ dari $N$, jika
	$\sum_{i \in I}N_i=N$, maka terdapat himpunan bagian hingga
	$I_0\subset I$ sedemikian sehingga $\sum_{i \in I_0}N_i=N$.
\end{lemma}
% segment-id: o014.aljabr2.chapter7.morita-theory-c.pf001
\begin{proof}
	Untuk arah ``hanya jika'', misalkan $x_1,\ldots,x_n$ suatu keluarga
	pembangkit $N$. Setiap $x_j$ termuat dalam suatu
	$\sum_{i\in I_j}N_i$, dengan $I_j\subset I$ hingga; cukup ambil
	$I_0=\bigcup_{j=1}^n I_j$. Untuk arah ``jika'', tinjau
	$N=\sum_{x\in N}xR$.
\end{proof}

% segment-id: o014.aljabr2.chapter7.morita-theory-c.p001
Sekarang kita kembali ke ekuivalensi Morita yang diperkenalkan dalam
Definisi--Proposisi~\ref{def:Morita-equiv}. Kita hanya membahas kasus modul
kanan.

% segment-id: o014.aljabr2.chapter7.morita-theory-c.t001
\begin{theorem}[Kiiti Morita]\label{prop:Morita-refined}
	Tinjau aljabar-$\Bbbk$ $A$ dan $B$. Misalkan
	$\mathrm{Equiv}(\cated{Mod}A,\cated{Mod}B)$ adalah kategori yang
	objek-objeknya semua ekuivalensi
	$F:\cated{Mod}A\to\cated{Mod}B$ dan morfisme-morfismenya isomorfisme antara
	ekuivalensi tersebut. Di sisi lain, definisikan kategori
	$\mathcal{P}(A,B)$ yang objek-objeknya adalah bimodul-$(A,B)$ $P$ yang
	memenuhi syarat-syarat berikut:
	\begin{itemize}
		\item sebagai modul-$B$ kanan, $P$ merupakan pembangkit projektif yang
		dibangkitkan secara hingga dari $\cated{Mod}B$;
		\item perkalian kiri menginduksi isomorfisme
		$A\rightiso\End_{\cated{Mod}B}(P)$.
	\end{itemize}
	Morfisme-morfismenya didefinisikan sebagai isomorfisme bimodul. Kita
	mempunyai funktor-funktor yang saling kuasi-invers berikut:
	\[
	\begin{tikzcd}[row sep=tiny]
		\mathrm{Equiv}(\cated{Mod}A, \cated{Mod}B) \arrow[shift left, r] & \mathcal{P}(A, B) \arrow[shift left, l] \\
		F \arrow[mapsto, r] & F(A) \\
		(\cdot) \dotimes{A} P & P \arrow[mapsto, l] .
	\end{tikzcd}
	\]
\end{theorem}
% segment-id: o014.aljabr2.chapter7.morita-theory-c.pf002
\begin{proof}
	Pertama, kita tunjukkan bahwa $F\mapsto F(A)$ terdefinisi dengan baik.
	Perhatikan bahwa $A$ merupakan pembangkit projektif yang dibangkitkan
	secara hingga dari $\cated{Mod}A$. Karena sifat menjadi pembangkit
	projektif maupun sifat dibangkitkan secara hingga sama-sama mempunyai
	karakterisasi kategoris (Lema~\ref{prop:fg-module-cat}), hal yang sama
	berlaku bagi $F(A)$. Selain itu, berdasarkan pertimbangan abstrak,
	% segment-id: o014.aljabr2.chapter7.morita-theory-c.q001
	\[
	F:\End_{\cated{Mod}A}(A)\rightiso
	\End_{\cated{Mod}B}(F(A)).
	\]
	Mudah diperiksa bahwa ini tepat merupakan homomorfisme
	$A\to\End_{\cated{Mod}B}(F(A))$ yang diberikan oleh perkalian kiri $A$
	pada $F(A)$.

	Jika funktor dari kanan ke kiri juga terdefinisi dengan baik, maka
	Teorema~\ref{prop:Morita} menyiratkan bahwa kedua funktor itu saling
	kuasi-invers. Dengan demikian, cukup dibuktikan pernyataan berikut: jika
	$P\in\Obj(\mathcal{P}(A,B))$, maka
	$F:=(\cdot)\dotimes{A}P$ merupakan ekuivalensi kategori.

	Ambil $P$ seperti di atas. Dalam \eqref{eqn:AB-adjunction-2} telah
	ditunjukkan bahwa adjoin kanan dari $(\cdot)\dotimes{A}P$ ialah
	$(\cdot)\dotimes{B}P^\vee$, dengan
	$P^\vee:=\Hom_{\cated{Mod}B}(P,B)$. Tujuannya adalah menunjukkan bahwa
	keduanya saling kuasi-invers. Dengan mengambil $Q=P^\vee$ dalam
	Definisi--Proposisi~\ref{def:Morita-equiv} (ii), selanjutnya cukup dibuktikan
	bahwa homomorfisme bimodul dalam
	Definisi~\ref{def:bimodule-ev-coev},
	% segment-id: o014.aljabr2.chapter7.morita-theory-c.q002
	\[
	\mathrm{ev}:P^\vee\dotimes{A}P\to B,\qquad
	\mathrm{coev}:A\to P\dotimes{B}P^\vee
	\simeq\End_{\cated{Mod}B}(P),
	\]
	keduanya merupakan isomorfisme.

	Pertama, tinjau $\mathrm{ev}$. Karena $P$ merupakan pembangkit dari
	$\cated{Mod}B$, terdapat homomorfisme surjektif
	$P^{\oplus I}\twoheadrightarrow B$. Khususnya, terdapat
	$q_1,\ldots,q_n\in P^\vee$ dan $p_1,\ldots,p_n\in P$ sedemikian sehingga
	$\sum_{i=1}^n q_i(p_i)=1_B$, atau dengan kata lain
	$\mathrm{ev}(\sum_i q_i\otimes p_i)=1_B$. Jadi, surjektivitas telah
	terbukti.

	Selanjutnya, definisikan operasi $P\times P^\vee\to A$ yang ditulis
	sebagai perkalian. Operasi ini dicirikan oleh kesamaan menyerupai
	asosiativitas berikut di dalam $P$, untuk semua $p,p'\in P$ dan
	$q\in P^\vee$:
	% segment-id: o014.aljabr2.chapter7.morita-theory-c.q003
	\[
	\underbracket{(pq)}_{\in A}p'
	= p\underbracket{(qp')}_{\in B}.
	\]
	Di sini $qp':=q(p')$. Memang, untuk $p$ dan $q$ tetap, ruas kanan
	merupakan unsur $\End_{\cated{Mod}B}(P)$, sedangkan perkalian kiri
	menginduksi $A\rightiso\End_{\cated{Mod}B}(P)$; dengan demikian,
	$pq\in A$ terdefinisi secara tunggal. Dari kesamaan di atas diperoleh
	pula bentuk asosiativitas kedua, yakni kesamaan berikut dalam $P^\vee$:
	% segment-id: o014.aljabr2.chapter7.morita-theory-c.q004
	\[
	\underbracket{(qp)}_{\in B}q'
	= q\underbracket{(pq')}_{\in A}.
	\]
	Cukup memeriksa bahwa kedua ruas memberikan nilai yang sama pada setiap
	$p'\in P$; verifikasi langsung ini diserahkan kepada pembaca.

	Sekarang kita dapat membuktikan injektivitas $\mathrm{ev}$. Pilih $p_i$
	dan $q_i$ seperti pada langkah sebelumnya. Jika
	$\mathrm{ev}(\sum_{j=1}^m q'_j\otimes p'_j)=0$, maka kedua bentuk
	asosiativitas di atas memberikan
	% segment-id: o014.aljabr2.chapter7.morita-theory-c.q005
	\begin{multline*}
		\sum_j q'_j \otimes p'_j
		= \sum_{i,j}(q_i p_i)q'_j\otimes p'_j
		= \sum_{i,j}q_i\underbracket{(p_iq'_j)}_{\in A}\otimes p'_j \\
		= \sum_{i,j}q_i\otimes(p_iq'_j)p'_j
		= \sum_{i,j}q_i\otimes p_i\underbracket{(q'_jp'_j)}_{\in B}
		= \sum_i q_i\otimes p_i\sum_j q'_jp'_j=0.
	\end{multline*}

	Terakhir, pernyataan bahwa $\mathrm{coev}$ merupakan isomorfisme tidak
	lain adalah bagian kedua dari definisi $\mathcal{P}(A,B)$.
\end{proof}

% segment-id: o014.aljabr2.chapter7.morita-theory-c.p002
Sebagai akibatnya, $A$ dan $B$ ekuivalen Morita jika dan hanya jika
$\mathcal{P}(A,B)$ tidak kosong, sedangkan pengkajian kategori
$\mathcal{P}(A,B)$ merupakan masalah teori modul yang sepenuhnya konkret.
Sebaliknya, kita juga dapat menggunakan ekuivalensi Morita untuk menulis
ulang definisi $\mathcal{P}(A,B)$ dalam bentuk yang lebih simetris.

% segment-id: o014.aljabr2.chapter7.morita-theory-c.c001
\begin{corollary}\label{prop:Morita-refined-symm}
	Untuk bimodul-$(A,B)$ $P$, modul tersebut merupakan objek kategori
	$\mathcal{P}(A,B)$ dalam Teorema~\ref{prop:Morita-refined} jika dan hanya
	jika syarat-syarat berikut berlaku:
	\begin{itemize}
		\item sebagai modul-$A$ kiri, $P$ merupakan pembangkit projektif yang
		dibangkitkan secara hingga dari $A\dcate{Mod}$, dan sebagai modul-$B$
		kanan, $P$ merupakan pembangkit projektif yang dibangkitkan secara
		hingga dari $\cated{Mod}B$;
		\item perkalian kiri menginduksi isomorfisme
		$A\rightiso\End_{\cated{Mod}B}(P)$, sedangkan perkalian kanan
		menginduksi isomorfisme
		$B^{\opp}\rightiso\End_{A\dcate{Mod}}(P)$.
	\end{itemize}
\end{corollary}
% segment-id: o014.aljabr2.chapter7.morita-theory-c.pf003
\begin{proof}
	Cukup dibuktikan arah ``hanya jika''. Misalkan
	$P\in\Obj(\mathcal{P}(A,B))$. Isomorfisme bimodul yang telah
	dibuktikan,
	% segment-id: o014.aljabr2.chapter7.morita-theory-c.q006
	\[
	P\dotimes{B}P^\vee\simeq A,\qquad
	P^\vee\dotimes{A}P\simeq B,
	\]
	tidak hanya menunjukkan bahwa
	$(\cdot)\dotimes{A}P:\cated{Mod}A\to\cated{Mod}B$ merupakan ekuivalensi,
	tetapi juga menyiratkan bahwa
	$P\dotimes{B}(\cdot):B\dcate{Mod}\to A\dcate{Mod}$ merupakan
	ekuivalensi. Oleh karena itu, versi modul kiri dari
	Teorema~\ref{prop:Morita-refined} menunjukkan bahwa $P$, sebagai
	modul-$A$ kiri, merupakan pembangkit projektif yang dibangkitkan secara
	hingga, dan bahwa perkalian kanan menginduksi
	$B^{\opp}\rightiso\End_{A\dcate{Mod}}(P)$. Sesungguhnya, cukup mengulangi
	paragraf pertama bukti Teorema~\ref{prop:Morita-refined} untuk modul kiri.
\end{proof}
